\section{A 10+2d Theory?}\label{sec:12d-supergravity-attempts}
There exists an algebra that admits a 10+2d interpretation and may unify the algebras of the various theories with 32 supercharges. 
Historically, this motivated a series of investigations into formulating a supersymmetric 10+2d theory of gravity.
These constructions will be reviewed for completeness.
The objective here is to present these attempts in an organised fashion, discuss their obstructions, and set the stage for the brane and effective action discussions later.

In 12d the complex Dirac spinor has 64 components and decomposes into two 32-component Weyl spinors.
In general, they are complex, but for the Lorentz group SO(10, 2),\footnote{We write SO(s, t) for the orthogonal group of a metric with $s$ spacelike and $t$ timelike directions, so SO(10, 2) has signature $(10, 2)$.} because
\begin{equation}
s-t\equiv 0 \pmod 8,
\end{equation}
one can impose a Majorana condition compatible with chirality to find Majorana-Weyl spinors with 32 real components \cite{VanProeyen:1999ni}.
The corresponding superalgebra contains a 2-form and a self-dual 6-form charge
\cite{vanHolten:1982mx}
\begin{equation}
\operatorname{Sym}^2[000001]_{SO(10, 2)} = [010000]_{SO(10, 2)} + [000020]_{SO(10, 2)}.
\label{OSp algebra}
\end{equation}
This algebra is often referred to as the $\mathrm{OSp}(1|32)$ algebra \cite{Bergshoeff:2000qu} or the F-theory algebra \cite{Hewson:1999tf}.
It reduces to the 11d algebra
\begin{equation}
\operatorname{Sym}^2[00001]_{SO(10, 1)} = [10000]_{SO(10, 1)} + [01000]_{SO(10, 1)} + [00002]_{SO(10, 1)}.
\end{equation}
Besides the 11d algebra, \eqref{OSp algebra} also reduces to the 10d type IIA and type IIB algebras \cite{Bergshoeff:2000qu,Bergshoeff:2000vg,Bergshoeff:2000vh}.
This unifying role made it a natural starting point for attempts to formulate a supersymmetric 10+2d theory of gravity \cite{Castellani:1982ke}, and for studying the objects such a theory would contain. 
For instance, 2+2d branes can propagate in (10, 2) spacetime \cite{Blencowe:1988sk}. 
Their BPS states \cite{Ueno:1999xa} and the fractions of supersymmetry they preserve \cite{Manvelyan:1998kx} have been classified. 
At a purely algebraic level, these ideas extend to proposals for unifying the various dualities in 12 and 13 dimensions \cite{Bars:1996dz,Bars:1996us,Bars:1996ab,Bars:1996hr}, and even in 14 dimensions \cite{Rudychev:1997ui,Bars:1997ug}.

However, the possibility of a 10+2d theory hinges on there being two time-like directions, and there are fundamental issues with formulating a supersymmetric theory with two times.
The little group of an SO(10, 2) theory is the non-compact SO(9, 1), with no nontrivial unitary irreducible representations \cite{Weinberg:1995mt}.
For unitary representations of the supersymmetry algebra, we may write the supercharge anticommutator 
\begin{equation}
\{Q_\alpha,Q_\beta^\dagger\}=\sum_n(\Gamma_n)_{\alpha\beta} Z_n,
\end{equation}
where $\Gamma_n$ are the antisymmetrised gamma-matrix structures and $Z_n$ the central charges. 
Because the RHS is symmetric and positive definite, we can diagonalise it.
Then we obtain fermionic raising and lowering operators, which implies that there are 256 states in the system. 
However, this cannot be carried out if the representation is non-unitary.
One might say, in any known one-time supergravity theory with 32 supercharges, there are 128+128 states, so start with this many as well. 
But the gravitino representation of SO(10) already has 144 states, exceeding the budget for fermions. There indeed is a way to add up to 144 + 144 states for SO(10) without too many scalars: a single gravitino in the fermion sector plus a graviton and two 2-forms in the bosonic sector. 
But this does not seem related to any known 10d or 11d multiplets.

Another issue that accompanies a non-compact little group is negative-norm states, or ``ghosts''.
Partially motivated by studying the 10+2d theory, a practical two-time framework has been developed \cite{Bars:1998cs,Bars:1998ui,Bars:1999nk,Bars:1999uq,Bars:2000qm}. 
The key ingredient is a local Sp(2, $\mathbb{R}$) gauge symmetry acting on the phase-space variables, treated as a two-component vector.
Different gauge choices lead to different one-time systems that share a common (d-2, 2) parent description.
The various approaches to formulating 10+2d theories largely follow this logic.

There have been some attempts at 10+2d SYM \cite{Nishino:1996wp,Sezgin:1997gr,Nishino:1997hk,Nishino:1997ia,Nishino:1999ck} and 10+2d supergravity \cite{Nishino:1997gq,Nishino:1997sw,Nishino:1998qn}, but none of them have achieved satisfactory results, as they either introduce null projectors that explicitly break SO(10, 2),
or fail to construct vielbeins due to the lack of a momentum generator \cite{Hewson:1999tf} in the algebra \eqref{OSp algebra}. 

After it was shown that the (2+2)-brane can propagate in 10+2d spacetime \cite{Blencowe:1988sk}, the $\mathcal{N}=(2,1)$ string was also studied \cite{Kutasov:1996fp,Kutasov:1996zm,Martinec:1996wn}. The $\mathcal{N}=1$ sector contains an effective 10+2d target space, with the $\mathcal{N}=2$ sector on a 2+2d target.
Later, an $\mathcal{N}=1$ superstring with a 2+2d target space was constructed \cite{Khviengia:1995pm,Lu:1995pn}.
Further work investigated whether self-dual gravity in 2+2 dimensions admits a stringy description \cite{Sezgin:1995fj} and related these self-dual 2+2d strings to a supersymmetric membrane action with $\mathrm{OSp}(1|32)$ and $\mathrm{OSp}(8|2)$ subgroup structure \cite{Ketov:1996gr}.
In parallel, a Green-Schwarz-type super 2+2d brane embedded in a 10+2d background was constructed \cite{Hewson:1996yh,Hewson:1997wv,Hewson:1998sw}.

There were also investigations into higher-dimensional bosonic field theories that, upon compactification and truncation, may produce the known theories.
In \cite{Khviengia:1997rh}, a 12d action with imaginary dilaton couplings was suggested as
\begin{equation}
2\kappa_{12}^2S
=\int d^{12}x\sqrt{-\widehat{g}}
\left[
\widehat R-\frac{1}{2}(\partial\widehat{\Phi})^2
-\frac{1}{2}e^{\frac{2i}{\sqrt{5}}\widehat{\Phi}}|\widehat{F}_5|^2
-\frac{1}{2}e^{\frac{i}{\sqrt{5}}\widehat{\Phi}}|\widehat{F}_4|^2
\right]
+\frac{\sqrt{3}}{4}\int \widehat{C}_4\wedge \widehat{F}_4 \wedge \widehat{F}_4
\label{Khviengia action}
\end{equation}
where $\widehat{F}_5 = d\widehat{C}_4$, $\widehat{F}_4 = d\widehat{C}_3$, $[\kappa_{12}^2]=L^{10}$, and $\widehat{\Phi}$ here is a 12d dilaton-like scalar.
We will use hats $\widehat{\phantom{x}}$ for higher-dimensional fields and straight letters $g_{mn},C_n,F_{n+1}$ for lower-dimensional fields. 
This will be discussed more clearly in Section~\ref{sec:covariant-unification}.
This action \eqref{Khviengia action} was subsequently studied by \cite{Bakhtiarizadeh:2018upg} at higher derivatives.
The most prominent feature of the proposal \eqref{Khviengia action} by \cite{Khviengia:1997rh} is the 12d dilaton and its imaginary couplings.
It was introduced based on a certain scalar invariant \cite{Lu:1995cs} across 11d and 10d type IIB brane solutions \cite{Khviengia:1997rh}.

More recently, it has been claimed that F-theory \cite{Vafa:1996xn} admits a 12d action \cite{Choi:2014vya,Choi:2015gia}, taking the form
\begin{equation}
2\kappa_{12}^2S = \int d^{11}x dy \sqrt{-\widehat{g}}
\left(
\widehat{R}-\frac{1}{2}|\widehat{F}_5|^2
\right)
+\frac{1}{6}
\int \widehat{C}_4\wedge \widehat{F}_4 \wedge \widehat{F}_4.
\label{Choi action}
\end{equation}
Here $\widehat{g}$ is the 12d metric, $\widehat{F}_5$ is a 5-form field strength.
All dependence on the twelfth coordinate $y$ is packaged into a scalar $r\equiv r(x^m,y)$, which also enters the metric.
In particular
\begin{equation}
\begin{gathered}
\widehat{C}_4(x,y) \equiv r(x,y) C_3\wedge dy,\quad
\widehat{F}_5(x,y) \equiv r(x,y) d C_3\wedge dy,\\
\widehat{g}_{mn} = g_{mn},\quad
\widehat{g}_{my} = 0,\quad
\widehat{g}_{yy} = r^2.
\label{eq:choi-embedding}
\end{gathered}
\end{equation}
Upon closer examination, we find \eqref{Choi action} is just 11d supergravity integrated over a spectator dimension.
This can be seen by substituting \eqref{eq:choi-embedding} into \eqref{Choi action} and evaluating $\int d^{11}x dy \sqrt{-\widehat{g}} \widehat{R}$.
The action \eqref{Choi action} reduces to
\begin{equation}
2\kappa_{12}^2 S =
\int d^{11}x \lambda(x)\sqrt{-g}
\left(R-\frac{1}{2}|F_4|^2\right)
+\frac{1}{6}\int \lambda(x) C_3\wedge F_4\wedge F_4
+\text{boundary},\quad
\lambda(x)\equiv \int dy r(x,y).
\end{equation}
This is the bosonic action of 11d supergravity, multiplied by an auxiliary field $\lambda(x)$ whose equation of motion imposes that the 11d bosonic Lagrangian vanishes.

In an attempt to unify the various string theory dualities, an \SLR{}$\times \mathbb{R}^+$ Exceptional Field Theory was proposed \cite{Berman:2015rcc,Rudolph:2017vzy}.
The idea is to realise unified dualities with extended coordinates, but this is not a 12d theory.

Related insights on a 12d viewpoint include proposals in which the type IIB theory is an edge mode of a gapped 12d bulk theory, motivated by a 6d generalisation of the fractional quantum Hall effect \cite{Heckman:2017uxe}, as well as the perspectives that 2+2 signature spacetime may be a tool for studying Lorentzian QFT \cite{Heckman:2018mxl,Heckman:2022peq}. 
We also note a recent supergravity construction of a locally supersymmetric $SO(10, 2)$-invariant action in $D=10+2$,
of MacDowell-Mansouri type and containing the Einstein-Hilbert term \cite{Castellani:2017vbi}, whose physical degrees of freedom and relation to standard 10d/11d supergravity remain to be clarified.

The persistent obstructions reviewed above (the non-compact little group, the absence of a momentum generator in the $\mathrm{OSp}(1|32)$ algebra, and the failure to match known 10d or 11d multiplets) suggest that a dynamical 10+2d supergravity with 32 supercharges reducing to the type IIB theory is not available.
A more modest strategy is then to ask whether specific solutions of the known 10d theories admit a 12d geometric interpretation, without requiring a full dynamical 12d theory.
This is the approach taken in the following sections.
